
\section{Requisitos del sistema}

\subsection{Introducción}

En esta sección se presentan los requisitos del Sistema de Registro de Eventos Interculturales (SREI). Su definición parte de las necesidades identificadas durante el levantamiento de información y de la especificación de los módulos de Gestión de Usuarios y Roles y Propuesta del Evento.

Para organizar los requisitos se consideran las categorías de negocio (BRS), interesados (StRS), sistema (SyRS), calidad (QR) y restricciones (C), siguiendo la estructura de la norma ISO/IEC/IEEE 29148:2018.

\subsection{Fuentes y técnica de elicitación}

La entrevista semiestructurada fue la técnica utilizada para conocer las necesidades de las personas relacionadas con la organización y participación en eventos interculturales. Permitió orientar las preguntas según el perfil de cada interesado y profundizar en los aspectos que requerían mayor explicación.

Como fuentes documentales se consideran el informe de entrevistas del proyecto y las especificaciones de requisitos de Gestión de Usuarios y Roles y Propuesta del Evento. También se dispone de grabaciones de audio correspondientes a entrevistas y validaciones. Estas grabaciones se conservan como evidencias; no se presentan como transcripciones.

La información recopilada se organizó de acuerdo con las necesidades de los interesados y se relacionó con los requisitos correspondientes. Los identificadores de las especificaciones originales se mantienen para facilitar su seguimiento.

\subsection{Perfiles de los interesados}

Los perfiles que participan en la definición y validación de los requisitos se presentan a continuación.

\begin{longtable}{p{0.29\textwidth} p{0.63\textwidth}}
\caption{Perfiles de los interesados del SREI}\\
\toprule
\textbf{Interesado} & \textbf{Rol y responsabilidad}\\
\midrule
\endfirsthead
\toprule
\textbf{Interesado} & \textbf{Rol y responsabilidad}\\
\midrule
\endhead
Ing. Jéssica Alexandra Ponce Ordoñez &
Product Owner. Coordinadora de la carrera de Software de la FCCDD, UTEQ. Proporciona requisitos.\\
Ing. Washington Alberto Chiriboga Casanova &
Product Owner. Decano de la Facultad de Ciencias de la Computación y Diseño Digital, UTEQ. Proporciona requisitos.\\
Msc. Elizabeth Johanna Contreras Jara &
Product Owner. Docente de Educación Inicial de la Facultad de Ciencias de la Educación, UTEQ. Proporciona requisitos.\\
Dr. Gleiston Cicerón Guerrero Ulloa &
Scrum Master. Facilitador externo y tutor académico del proceso Scrum.\\
Washington Germán Apunte Pazmiño &
Integrante del equipo de desarrollo.\\
Jordan Josue Loor Suárez &
Integrante del equipo de desarrollo.\\
Steffanny del Rocío Loor Suárez &
Integrante del equipo de desarrollo.\\
Estudiante participante &
Usuario final que consulta los eventos, se registra y confirma su asistencia.\\
\bottomrule
\end{longtable}

\subsection{Clasificación de los requisitos}

Los requisitos se organizan en cinco grupos:

\begin{itemize}
    \item \textbf{BRS:} requisitos de negocio que expresan los objetivos institucionales.
    \item \textbf{StRS:} necesidades y expectativas de los interesados.
    \item \textbf{SyRS:} requisitos funcionales que describen las operaciones que debe realizar el sistema.
    \item \textbf{QR:} requisitos relacionados con la calidad, como usabilidad, rendimiento, seguridad y compatibilidad.
    \item \textbf{C:} restricciones que delimitan las tecnologías, formatos y condiciones de funcionamiento.
\end{itemize}

\subsection{Requisitos de negocio}

La especificación contempla dos requisitos de negocio principales:

\begin{itemize}
    \item \textbf{BRS-01:} controlar de forma centralizada el acceso al módulo y las acciones permitidas según el rol de cada persona.
    \item \textbf{BRS-02:} formalizar las propuestas de eventos interculturales, de manera que los eventos publicados cuenten con un informe PDF de respaldo y un trámite verificable.
\end{itemize}

\subsection{Requisitos de los interesados}

Se identifican las siguientes necesidades:

\begin{itemize}
    \item \textbf{StRS-01:} el Administrador necesita crear, consultar, actualizar y dar de baja a los usuarios.
    \item \textbf{StRS-02:} el Administrador necesita definir los roles y sus respectivos permisos.
    \item \textbf{StRS-03:} el Docente necesita registrar una propuesta de evento, adjuntar el informe PDF y enviarla a revisión.
    \item \textbf{StRS-04:} el Docente necesita registrar los campos particulares de cada tipo de evento, sin tener que utilizar campos genéricos.
\end{itemize}

\subsection{Historias de usuario}

A partir de los requisitos funcionales especificados para los módulos de Gestión de Usuarios y Roles y Propuesta del Evento, se presentan diez historias de usuario. Cada una incluye criterios de aceptación, prioridad y trazabilidad hacia los requisitos originales.

\begin{longtable}{p{0.07\textwidth} p{0.33\textwidth} p{0.36\textwidth} p{0.09\textwidth} p{0.11\textwidth}}
\caption{Historias de usuario y criterios de aceptación}\\
\toprule
\textbf{ID} & \textbf{Historia de usuario} & \textbf{Criterios de aceptación} & \textbf{Prioridad} & \textbf{Trazabilidad}\\
\midrule
\endfirsthead
\toprule
\textbf{ID} & \textbf{Historia de usuario} & \textbf{Criterios de aceptación} & \textbf{Prioridad} & \textbf{Trazabilidad}\\
\midrule
\endhead

HU-01 &
Como Administrador, quiero registrar usuarios con sus datos personales y correo institucional para mantener organizada la información de las personas que utilizan el sistema. &
El sistema solicita cédula, nombres, apellidos, correo institucional, contraseña y rol. El correo debe pertenecer al dominio @uteq.edu.ec. La cédula debe contener diez dígitos numéricos. &
Alta &
RF-01, RF-02, RF-03\\

HU-02 &
Como Administrador, quiero evitar registros duplicados para mantener información única de cada usuario. &
El sistema verifica que la cédula y el correo no estén registrados previamente. Si alguno está duplicado, rechaza el registro e informa el motivo. &
Alta &
RF-04\\

HU-03 &
Como Administrador, quiero modificar y desactivar usuarios para mantener actualizada la información y controlar quién puede acceder al sistema. &
El sistema permite modificar los datos del usuario y cambiar su estado. La desactivación conserva el registro y su historial de acciones. &
Alta &
RF-06, RF-07\\

HU-04 &
Como Administrador, quiero consultar y filtrar los usuarios registrados para encontrar rápidamente una cuenta. &
El listado permite filtrar por rol, estado, nombre, apellido o cédula. &
Media &
RF-08\\

HU-05 &
Como Administrador, quiero administrar los roles y permisos para controlar las acciones que puede realizar cada tipo de usuario. &
El sistema permite crear, consultar, modificar y eliminar roles y permisos. También permite asignar y revocar permisos, sin duplicar asignaciones. &
Alta &
RF-09, RF-10, RF-11, RF-12, RF-13, RF-14\\

HU-06 &
Como Docente, quiero registrar una propuesta de evento con sus datos principales para enviarla a revisión. &
El formulario solicita lugar, tipo de evento, título, descripción, fecha, hora de inicio y hora de fin. La hora de finalización debe ser posterior a la de inicio. &
Alta &
RF-17, RF-18, RF-19, RF-20, RF-21\\

HU-07 &
Como Docente, quiero adjuntar una imagen al evento y declarar su origen para identificar si fue cargada manualmente o generada mediante inteligencia artificial. &
El sistema permite seleccionar el origen SUBIDA o IA. Si se selecciona IA, exige el prompt. Si se selecciona SUBIDA, el campo del prompt debe permanecer vacío. &
Media &
RF-22, RF-23, RF-24\\

HU-08 &
Como Docente, quiero agregar información específica según el tipo de evento para registrar correctamente sus características. &
El sistema permite registrar pares de nombre y valor para los campos adicionales. No permite dos campos con el mismo nombre dentro de un evento. &
Media &
RF-25, RF-26\\

HU-09 &
Como Docente, quiero adjuntar el informe PDF y enviar la propuesta para que la Coordinación pueda revisarla. &
El informe PDF es obligatorio. Al enviar la propuesta, el sistema registra el evento y su propuesta asociada, asigna el estado PENDIENTE y genera un código QR. &
Alta &
RF-27, RF-28, RF-29, RF-30, RF-31\\

HU-10 &
Como Coordinador, quiero consultar el informe de respaldo y revisar las propuestas enviadas para controlar el proceso de aprobación de los eventos. &
El sistema restringe la visualización del informe PDF al rol COORDINADOR.  &
Alta &El registro de propuestas se limita al rol DOCENTE con el permiso EVENTO\_PROPONER.
RF-32, RF-33\\

\bottomrule
\end{longtable}

\subsection{Requisitos no funcionales}

Los requisitos no funcionales establecen condiciones de calidad que deben considerarse durante el desarrollo y funcionamiento del sistema. Se utilizan los identificadores originales de la especificación y se relacionan con características de calidad de ISO/IEC 25010.

\begin{longtable}{p{0.11\textwidth} p{0.28\textwidth} p{0.24\textwidth} p{0.25\textwidth}}
\caption{Requisitos no funcionales y criterios de verificación}\\
\toprule
\textbf{ID} & \textbf{Requisito} & \textbf{Criterio verificable} & \textbf{Característica ISO/IEC 25010}\\
\midrule
\endfirsthead
\toprule
\textbf{ID} & \textbf{Requisito} & \textbf{Criterio verificable} & \textbf{Característica ISO/IEC 25010}\\
\midrule
\endhead

PERS-01 &
El registro de usuarios y la actualización del catálogo de roles y permisos deben completarse en un tiempo no mayor a 3 segundos, bajo condiciones normales de operación. &
Medir desde el envío del formulario hasta la confirmación visible. El tiempo no debe superar los 3 segundos. &
Eficiencia del desempeño: comportamiento temporal.\\

PERS-02 &
El registro de una propuesta, incluyendo el informe PDF y la imagen, debe completarse en un tiempo no mayor a 5 segundos para archivos de hasta 10 MB, bajo condiciones normales de operación. &
Realizar el registro con archivos dentro del límite indicado y medir el tiempo de respuesta. &
Eficiencia del desempeño: comportamiento temporal y capacidad.\\

USARS-01 &
El formulario de registro de usuario debe mostrar mensajes específicos junto a los campos que presenten errores de validación. &
Comprobar los tres casos indicados: correo fuera del dominio institucional, cédula incorrecta y datos duplicados. &
Capacidad de interacción: asistencia al usuario.\\

USARS-02 &
El formulario de eventos debe mostrar u ocultar dinámicamente el campo de aforo según el tipo de espacio seleccionado. &
Verificar los dos estados: espacio CERRADO, que exige aforo mayor a cero, y espacio ABIERTO, que mantiene el campo sin valor. &
Capacidad de interacción: operabilidad.\\

USARS-03 &
El listado de usuarios debe ofrecer controles de búsqueda y filtrado sin recargar la página completa. &
Comprobar los tres criterios de filtrado: rol, estado y nombre o cédula. &
Capacidad de interacción: operabilidad.\\

COMPRS-02 &
El registro de usuarios y las propuestas de eventos deben funcionar en Google Chrome, Mozilla Firefox y Microsoft Edge, sin pérdida de funcionalidad. &
Comprobar las dos funcionalidades principales en los tres navegadores indicados. &
Compatibilidad: adaptabilidad y coexistencia.\\

\bottomrule
\end{longtable}

Además, se consideran los siguientes requisitos de seguridad establecidos en la especificación:

\begin{itemize}
    \item \textbf{SECRS-01:} las contraseñas deben almacenarse mediante un hash de una sola dirección, sin conservarlas ni transmitirlas en texto plano.
    \item \textbf{SECRS-02:} la autenticación debe utilizar tokens JWT y verificar el rol ADMINISTRADOR y el permiso USUARIO\_GESTIONAR en cada solicitud de gestión.
    \item \textbf{SECRS-03:} el registro de propuestas debe comprobar el rol DOCENTE y el permiso EVENTO\_PROPONER.
    \item \textbf{SECRS-04:} la visualización y descarga del informe PDF debe estar restringida al rol COORDINADOR.
\end{itemize}

Los criterios de verificación descritos permiten comprobar las condiciones indicadas en los requisitos. Su inclusión en este documento no significa que las pruebas de aceptación o rendimiento ya se hayan ejecutado.

\subsection{Restricciones del sistema}

La especificación establece las siguientes restricciones:

\begin{itemize}
    \item \textbf{C-01:} el backend debe implementarse con Java 17 y Spring Boot, mientras que el frontend debe utilizar Angular con TypeScript.
    \item \textbf{C-02:} el correo electrónico registrado debe pertenecer al dominio institucional @uteq.edu.ec.
    \item \textbf{C-03:} el informe adjunto a la propuesta debe entregarse exclusivamente en formato PDF.
    \item \textbf{C-04:} los roles permitidos son COORDINADOR, DOCENTE, ESTUDIANTE y ADMINISTRADOR.
\end{itemize}

\subsection{Trazabilidad de los requisitos}

La trazabilidad permite relacionar las necesidades identificadas con las historias de usuario y los requisitos funcionales que las respaldan. La siguiente tabla resume esta relación.

\begin{longtable}{p{0.17\textwidth} p{0.30\textwidth} p{0.20\textwidth} p{0.23\textwidth}}
\caption{Matriz resumida de trazabilidad}\\
\toprule
\textbf{Fuente} & \textbf{Necesidad} & \textbf{Historia} & \textbf{Requisitos relacionados}\\
\midrule
\endfirsthead
\toprule
\textbf{Fuente} & \textbf{Necesidad} & \textbf{Historia} & \textbf{Requisitos relacionados}\\
\midrule
\endhead

Administrador &
Gestionar cuentas de usuarios. &
HU-01, HU-02, HU-03, HU-04 &
RF-01 a RF-08\\

Administrador &
Administrar roles y permisos. &
HU-05 &
RF-09 a RF-16\\

Docente &
Registrar propuestas de eventos. &
HU-06 &
RF-17 a RF-21\\

Docente &
Asociar imágenes y registrar su origen. &
HU-07 &
RF-22 a RF-24\\

Docente &
Registrar campos particulares de cada evento. &
HU-08 &
RF-25, RF-26\\

Docente &
Adjuntar el informe y enviar la propuesta. &
HU-09 &
RF-27 a RF-31\\

Coordinador &
Consultar los informes de respaldo y controlar el acceso. &
HU-10 &
RF-32, RF-33\\

\bottomrule
\end{longtable}

\subsection{Consideraciones de validación}

Los requisitos presentados conservan los identificadores de la especificación documental utilizada como fuente. La validación técnica de cada funcionalidad deberá realizarse durante las pruebas del sistema, comprobando los criterios de aceptación y las restricciones correspondientes.

Las grabaciones de audio de las entrevistas y validaciones se conservan como evidencias del proceso de levantamiento. Cuando se incorporen observaciones adicionales de los interesados, estas deberán registrarse y relacionarse con el requisito afectado, sin modificar su trazabilidad original.
